% Options for packages loaded elsewhere
\PassOptionsToPackage{unicode}{hyperref}
\PassOptionsToPackage{hyphens}{url}
\PassOptionsToPackage{dvipsnames,svgnames,x11names}{xcolor}
%
\documentclass[
  10pt,
  a4paper,
]{article}
\usepackage{amsmath,amssymb}
\usepackage{iftex}
\ifPDFTeX
  \usepackage[T1]{fontenc}
  \usepackage[utf8]{inputenc}
  \usepackage{textcomp} % provide euro and other symbols
\else % if luatex or xetex
  \usepackage{unicode-math} % this also loads fontspec
  \defaultfontfeatures{Scale=MatchLowercase}
  \defaultfontfeatures[\rmfamily]{Ligatures=TeX,Scale=1}
\fi
\usepackage{lmodern}
\ifPDFTeX\else
  % xetex/luatex font selection
\fi
% Use upquote if available, for straight quotes in verbatim environments
\IfFileExists{upquote.sty}{\usepackage{upquote}}{}
\IfFileExists{microtype.sty}{% use microtype if available
  \usepackage[]{microtype}
  \UseMicrotypeSet[protrusion]{basicmath} % disable protrusion for tt fonts
}{}
\makeatletter
\@ifundefined{KOMAClassName}{% if non-KOMA class
  \IfFileExists{parskip.sty}{%
    \usepackage{parskip}
  }{% else
    \setlength{\parindent}{0pt}
    \setlength{\parskip}{6pt plus 2pt minus 1pt}}
}{% if KOMA class
  \KOMAoptions{parskip=half}}
\makeatother
\usepackage{xcolor}
\usepackage[margin=25mm]{geometry}
\setlength{\emergencystretch}{3em} % prevent overfull lines
\providecommand{\tightlist}{%
  \setlength{\itemsep}{0pt}\setlength{\parskip}{0pt}}
\setcounter{secnumdepth}{5}
\usepackage{microtype}
\usepackage{mathtools}
\usepackage{amssymb}
\usepackage{newtxtext}
\usepackage{newtxmath}
\usepackage{enumitem}
\setlist{nosep}
\ifLuaTeX
  \usepackage{selnolig}  % disable illegal ligatures
\fi
\usepackage{bookmark}
\IfFileExists{xurl.sty}{\usepackage{xurl}}{} % add URL line breaks if available
\urlstyle{same}
\hypersetup{
  pdftitle={Upper-half-plane radiative mode exclusion for subextremal Kerr--Newman black holes},
  colorlinks=true,
  linkcolor={Maroon},
  filecolor={Maroon},
  citecolor={Blue},
  urlcolor={Blue},
  pdfcreator={LaTeX via pandoc}}

\title{Upper-half-plane radiative mode exclusion for subextremal
Kerr--Newman black holes}
\author{}
\date{}

\begin{document}
\maketitle

\textbf{Abstract.} We prove upper-half-plane radiative mode exclusion
for the linearized Einstein--Maxwell equations on every strictly
subextremal, asymptotically flat Kerr--Newman black hole. For each fixed
azimuthal number and frequency with positive imaginary part, outgoing
modes satisfying the regular axis and future-horizon conditions have
vanishing gravitational and electromagnetic radiative data. For nonzero
charge, the proof treats the full coupled first-order radiative PDE
system, using complex angular generalized-root completeness, a
division-free radial reduction, a distributional horizon extension, and
a coercive transformed energy identity. The Kerr case is closed by a
full-field spectral reduction and a spin-reversal lemma that covers
exceptional horizon frequencies without assuming a nonzero leading
horizon coefficient. The result concerns radiative modes and does not
assert a decay theorem or a classification of nonradiative gauge and
parameter perturbations.

\section{Introduction and theorem}\label{introduction-and-theorem}

The separation of the coupled gravitational and electromagnetic
radiative perturbation equations on subextremal Kerr--Newman backgrounds
{[}H{]} makes it possible to study the non-real mode problem directly.
The real-frequency exclusion established in {[}H{]} does not by itself
exclude frequencies with positive imaginary part. Here we prove that
exclusion for the radiative Einstein--Maxwell variables in the complete
strictly subextremal parameter domain. The zero-charge Kerr case is
treated separately, using the negative-spin version of the classical
separated Kerr result {[}W, AMPW{]} and an explicit spin-reversal
argument that includes exceptional horizon frequencies. Rigorous linear
stability is known in the restricted weak-charge and slow-rotation
regime {[}He{]}; earlier numerical studies of the coupled system
{[}DGS{]} and calculations of separated Kerr--Newman quasinormal modes
{[}DCH{]} provide complementary evidence in broader parameter ranges.
Recent work on radiative amplitudes and algebraic separation {[}M, Z{]}
supplies further context. We use geometrized units throughout.

Fix real parameters \[
 M>0,\qquad a^2+Q^2<M^2,
\] an integer \(m\), and a frequency \(\omega\in\mathbb C\) with
\(\gamma:=\operatorname{Im}\omega>0\). Time and azimuthal dependence is
\(e^{-i\omega t+im\phi}\).

\textbf{Theorem (upper-half-plane radiative mode exclusion).} Every
outgoing smooth linearized Einstein--Maxwell perturbation on this
asymptotically flat Kerr--Newman exterior has vanishing gravitational
and electromagnetic radiative data.

For \(Q\ne0\), the stronger sufficient assertion proved below is that
both entire four-component radiative first-order PDE systems in Hintz
{[}H{]}, with their prescribed axis, future-horizon and infinity
conditions, have only the zero solution. Their equivalence to the
transformed equations used here is algebraic and valid for complex
frequency. For \(Q=0\), the Kerr section below combines the classical
separated Kerr mode-stability theorem in its negative-spin form {[}W,
AMPW{]} with a spin-reversal lemma and an explicit full-PDE coverage
argument.

Put \[
 \Delta=r^2-2Mr+a^2+Q^2,\quad r_\pm=M\pm\delta,\quad
 \delta=\sqrt{M^2-a^2-Q^2}>0,
\] \[
 K=-(r^2+a^2)\omega+am,\qquad
 H(r)=\exp\left(i\int_{2r_+}^r K(s)/\Delta(s)\,ds\right).
\] Future-horizon regularity means that \(H^{-1}\Delta^b\psi_b\),
\(H^{-1}\Delta^2k\), \(H^{-1}\Delta s\), \(H^{-1}\Delta^{-2}n\), and
\(H^{-1}\Delta^{-1}\ell\) are smooth up to \(r=r_+\), using the
variables of the relevant spin system. This is smoothness in the regular
ingoing tetrad. At infinity, \(H\) times each original radiative
variable has polynomial conormal bounds, meaning that every fixed
collection of \(r\partial_r\) and angular derivatives has a polynomial
bound. The physical outgoing branch decays exponentially because
\(\gamma>0\). Axis smoothness means smoothness of the relevant
spin-weighted sections.

The theorem concerns radiative mode stability. Classification of
nonradiating perturbations as gauge and parameter variations, decay,
nonlinear stability, extremal backgrounds, and bounds uniform to
extremality are outside this claim.

\section{Source system and
transformations}\label{source-system-and-transformations}

\subsection{Constants and the regular
PDE}\label{constants-and-the-regular-pde}

Assume \(Q\ne0\), and define \[
 d=M^2-Q^2,\quad \eta=\frac{2Q^2}{d},\quad
 h=\eta+2=\frac{2M^2}{d},\quad
 \kappa=\eta+1=\frac{M^2+Q^2}{d},
\] \[
 c=\frac{2MQ}{d},\qquad \beta=\frac{3M}{2Q},\qquad
 \xi_0=\eta\omega,\qquad \widetilde K=K-\xi_0\Delta .
\] Then \[
 \kappa^2-c^2=1,\qquad \beta c=\frac{3h}{2}.
 \tag{1}
\] Let \(\sigma\in\{+1,-1\}\) distinguish the positive and negative spin
systems. The exact transformed equation {[}H, equations (7)--(8){]} is
\[
 \sqrt\Delta\,\partial_r Z
 =(\mathcal A_\sigma(\omega)+R_0\mathcal I+R_1\mathcal J)Z,
 \tag{2}
\] where \[
 R_0=\frac{3\Delta'}{4\sqrt\Delta}
       -i\sigma\sqrt\Delta\left(\frac K\Delta-\eta\omega\right),\qquad
 R_1=-\sigma\frac{2\omega Q(Mr-Q^2)}d+i\beta.
 \tag{3}
\] The constant Hermitian matrices obey \[
 \mathcal I^2=\mathcal J^2=1,\quad
 \mathcal I\mathcal J=-\mathcal J\mathcal I,\quad
 \{\mathcal A,\mathcal I\}=\{\mathcal A,\mathcal J\}=0.
 \tag{4}
\] Thus \(\mathcal S=\mathcal I\mathcal J\) has eigenvalues \(i,-i\) and
commutes with \(\mathcal A\). On \(L^2((0,\pi),d\theta;\mathbb C^4)\),
\[
 \mathcal A_\sigma=
 \begin{pmatrix}i\partial_\theta&i\Theta\\-i\Theta&-i\partial_\theta\end{pmatrix},
 \qquad \Theta=\begin{pmatrix}W&T\\T&-W\end{pmatrix},
 \tag{5}
\] \[
 W=\frac32\cot\theta+\sigma m\csc\theta-\sigma a\kappa\omega\sin\theta,
 \qquad T=-\beta-\sigma ac\omega\cos\theta.
\] One may take
\(\mathcal I=\left(\begin{smallmatrix}0&I_2\\I_2&0\end{smallmatrix}\right)\)
and
\(\mathcal J=\left(\begin{smallmatrix}0&-J_2\\J_2&0\end{smallmatrix}\right)\),
where
\(J_2=\left(\begin{smallmatrix}0&-1\\1&0\end{smallmatrix}\right)\). All
relations are complex algebraic identities; no conjugation of \(\omega\)
is used.

\subsection{Invertibility and boundary
fidelity}\label{invertibility-and-boundary-fidelity}

Write \(\rho=r+ia\cos\theta\), \(\rho^*=r-ia\cos\theta\), and \[
 X=\begin{pmatrix}
 Q(\sqrt\Delta+a\sin\theta)/\rho&-i(M-Q^2/\rho)\\
 i(M-Q^2/\rho^*)&Q(\sqrt\Delta-a\sin\theta)/\rho^*
 \end{pmatrix}.
 \tag{6}
\] Its determinant is \(-d\ne0\). The source transformation is \[
 y_\pm^{(+)}=\Delta^{5/4}\sqrt{\sin\theta}
 \begin{pmatrix}\psi_2\pm i\Delta^{-1/2}\psi_1\\
 -2Q(k\pm i\Delta^{-1/2}s)\end{pmatrix},
\] \[
 y_\pm^{(-)}=\Delta^{-3/4}\sqrt{\sin\theta}
 \begin{pmatrix}\psi_{-2}\mp i\Delta^{1/2}\psi_{-1}\\
 2Q(n\mp i\Delta^{1/2}\ell)\end{pmatrix},
\] and \(Z=(z_+,z_-)^t\), where \(z_+=\operatorname{adj}(X)y_+\),
\(z_-=X^*y_-\). The matrix \(X\) depends on the real background and
coordinates, not frequency. All factors are invertible in the open
exterior. The substitutions in the original first-order equations give
(2) without dividing by any coefficient of a second-order elimination.

Let \(\Pi_\pm=(1\pm\mathcal I)/2\). The source transformation gives \[
\begin{array}{c|cc}
 &\Pi_+Z&\Pi_-Z\\ \hline
 \sigma=+1&H\Delta^{-1/4}C^\infty&H\Delta^{-3/4}C^\infty\\
 \sigma=-1&H\Delta^{3/4}C^\infty&H\Delta^{5/4}C^\infty .
\end{array}
 \tag{7}
\] Smoothness may be interpreted into the angular Hilbert space.
Physical axis regularity implies this condition and the maximal angular
domain at each interior radius. At infinity the transformed components
are \(H^{-1}\) times polynomially bounded functions; the remaining
transformation factors have polynomial bounds.

For fixed \(Q\ne0\), the phantom tetrad gauge is legitimate because the
background Maxwell component \(\Phi_1=Q/(2\rho^{*2})\) is nowhere zero
in the exterior. The radiative variables are the rescaled phantom-gauge
Weyl components \[
 \psi_2=\Psi_0^{(1)},\quad
 \psi_1=\sqrt2\rho^*\Psi_1^{(1)},\quad
 \psi_{-1}=\rho^{*3}\Psi_3^{(1)}/\sqrt2,\quad
 \psi_{-2}=\rho^{*4}\Psi_4^{(1)}.
 \tag{8}
\] The exact radiative data meant in the theorem are
\(\mathcal G_\pm,\mathcal E_\pm\) of {[}M, equations (5.1)--(5.2){]}: \[
\begin{aligned}
 \mathcal G_+&=\delta\Psi_0,&
 \mathcal G_-&=\rho^{*4}\delta\Psi_4,\\
 \mathcal E_+&=\frac{\rho^*}{3M}
     (Q\delta\Psi_1-3\rho^{*2}\Psi_2^{(0)}\delta\Phi_0),\\
 \mathcal E_-&=\frac{\rho^{*3}}{3M}
     (Q\delta\Psi_3-3\rho^{*2}\Psi_2^{(0)}\delta\Phi_2).
\end{aligned}
 \tag{8a}
\] The mixed combinations are invariant under the infinitesimal null
rotations used to impose phantom gauge. In that gauge, \[
 \mathcal G_\pm=\psi_{\pm2},\qquad
 \mathcal E_+=\frac{Q}{3\sqrt2M}\psi_1,\qquad
 \mathcal E_-=\frac{\sqrt2Q}{3M}\psi_{-1}.
 \tag{8b}
\] Thus proving the four \(\psi\)'s vanish proves vanishing of these
well-defined radiative data. At finite charge, \(\mathcal E_\pm\)
contain both curvature and Maxwell perturbations. Their zero-charge
limits in a nonsingular tetrad are \(\delta\Phi_0\) and
\(\rho^{*2}\delta\Phi_2\), respectively. These local algebraic
identifications do not require a real-frequency flux formula. We use the
one-way implication from physical perturbations to the necessary
first-order systems. Reconstruction of arbitrary subsystem solutions is
not an assumption.

\section{Complex angular analysis}\label{complex-angular-analysis}

\subsection{Fixed domain and exact zero-frequency
spectrum}\label{fixed-domain-and-exact-zero-frequency-spectrum}

On \(\ker(\mathcal S-i)\), a constant unitary identification gives \[
 B_\sigma(\omega)=B_0+\omega V,\qquad
 B_0=i\sigma_3\partial_\theta+g\sigma_2-\beta\sigma_1,\qquad
 g=\frac32\cot\theta+\sigma m\csc\theta,
 \tag{9}
\] \[
 V=-\sigma a(\kappa\sin\theta\,\sigma_2+c\cos\theta\,\sigma_1).
\] Here \(\sigma_j\) are the Hermitian Pauli matrices; the scalar system
sign \(\sigma\) is distinct.

The compactly supported smooth realization of \(B_0\) is essentially
self-adjoint, with maximal distributional \(L^2\) domain. This follows
also from {[}H, Lemma 1{]} at zero frequency. The endpoint indicial
roots are \(\pm(3\cos\theta_0/2+\sigma m)\), nonzero half-integers. The
root at or below \(-1/2\) is not square integrable. The deficiency
integration has zero boundary flux and forces every deficiency vector to
vanish. The bounded Hermitian multiplier \(V\) leaves the operator
domain fixed at complex \(\omega\).

A constant Hadamard rotation puts \(B_0\) in the form \[
 \begin{pmatrix}-\beta&i(\partial_\theta+g)\\
 i(\partial_\theta-g)&\beta\end{pmatrix}.
 \tag{10}
\] Its square minus \(\beta^2\) has diagonal entries \[
 L_s=-\partial_\theta^2+
 \frac{m^2+s^2-\tfrac14+2\sigma ms\cos\theta}{\sin^2\theta}
 -\frac94,\qquad s=1,2.
 \tag{11}
\] Let \(A_s=|\sigma m+s|\), \(B_s=|\sigma m-s|\). The functions \[
 \sqrt{\sin\theta}(1-\cos\theta)^{A_s/2}
 (1+\cos\theta)^{B_s/2}P_n^{(A_s,B_s)}(\cos\theta),\qquad n\ge0,
 \tag{12}
\] have eigenvalue \(\ell(\ell+1)-2\), where \(\ell=n+\max(|m|,s)\), by
the Jacobi differential equation. The change of variable
\(x=\cos\theta\) identifies their norms with Jacobi weighted norms, so
polynomial density proves completeness. At critical exponent \(1/2\),
the leading coefficient of a potentially nonintegrable first-order image
cancels. The displayed functions and their first-order images belong to
the maximal domain. Thus (11) is precisely the square of (10), with no
unauthorized logarithmic second-order boundary extension.

For \(\ell\ge\max(|m|,2)\), the off-diagonal part of (10) couples the
two one-dimensional eigenspaces in (12), giving the complete paired
spectrum \[
 \nu_{\ell,\pm}=\pm\sqrt{\beta^2+(\ell-1)(\ell+2)}.
 \tag{13}
\] If \(|m|\le1\), there is additionally one unpaired eigenvalue
\(\nu=-\beta\), with upper component proportional to
\((\sin\theta)^{3/2}\tan(\theta/2)^{\sigma m}\) and lower component
zero. There is no unpaired \(+\beta\) eigenvalue in this chiral sector.
Let \(P_0\) be that possible rank-one projection. In quadratic-form
sense on \(\operatorname{Dom}B_0\), \[
 T_0:=B_0^2-\beta^2\ge4(1-P_0)\ge0.
 \tag{14}
\]

\subsection{The required nonnegative angular
form}\label{the-required-nonnegative-angular-form}

For real \(x\le-h\), set \[
 K_\sigma(x)=B_0^2-\beta^2+\frac c x(B_0-\sigma\beta).
 \tag{15}
\] This form is nonnegative for both signs of \(Q\) and \(\sigma\).
Indeed, put \(q=|Q|/M\in(0,1)\), \(b=|\beta|=3/(2q)\). Then
\(|c/x|\le q\), with sign opposite to \(\beta\). For \(\sigma=-1\), the
smaller nonzero-branch value is bounded by \[
 v^2-b^2-q(v+b),\qquad v\ge\sqrt{b^2+4}.
\] It increases with \(v\), and its endpoint value is \[
 \epsilon_q=\frac{5-\sqrt{9+16q^2}}2
 =\frac{8(1-q^2)}{5+\sqrt{9+16q^2}}>0.
\] The unpaired branch has value zero. Hence \[
 K_-(x)\ge\epsilon_q(1-P_0).
 \tag{16}
\] For \(\sigma=+1\), the smaller nonzero-branch value is at least
\((11-\sqrt{9+16q^2})/2>3\), and the unpaired value is
\(-2\beta c/x=-3h/x>0\). Consequently \[
 K_+(x)\ge-\frac{3h}{x}P_0+
 \frac{11-\sqrt{9+16q^2}}2(1-P_0)\ge0.
 \tag{17}
\] No uniform estimate at the extremal boundary \(q=1\) is used.

For a normalized genuine eigenvector \(B_\sigma(\omega)u=\lambda u\),
introduce real numbers \[
\begin{gathered}
 L=\langle u,B_0u\rangle,\quad v=\langle u,Vu\rangle,\quad
 T=\|B_0u\|^2-\beta^2,\\
 C=2\operatorname{Re}\langle B_0u,Vu\rangle,\quad D=\|Vu\|^2.
\end{gathered}
\] With the inner product conjugate-linear in its first entry, \[
 \lambda=L+\omega v,\qquad
 \lambda^2=\beta^2+T+\omega C+\omega^2D .
 \tag{18}
\] The latter identity is valid using only the first-order domain: \[
 \langle B_\sigma(\omega)^*u,B_\sigma(\omega)u\rangle
 =\lambda\langle B_\sigma(\omega)^*u,u\rangle
 =\lambda\langle u,B_\sigma(\omega)u\rangle=\lambda^2.
\] Expanding the left side proves (18). This does not presume
self-adjointness at complex frequency or a square-operator domain. The
Pauli algebra gives \[
 V^2=a^2(c^2+\sin^2\theta)1,\qquad \|V\|=|a|\kappa.
 \tag{19}
\] The eigenvector in these expectations is fixed independently of the
transform coordinate \(x\).

\subsection{Generalized-root
completeness}\label{generalized-root-completeness}

For each fixed \(m\), (13) implies that one resolvent of \(B_0\), and of
the full \(\mathcal A_0\), belongs to Schatten class \(\mathcal S_p\)
for every \(p>1\), in particular \(\mathcal S_2\). The bounded
perturbation preserves this property at a sufficiently distant resolvent
point. If \(C_\omega=|\omega|\|\mathcal V\|\), where \(\mathcal V\) is
the full-space frequency multiplier, \[
 \|(\mathcal A(\omega)-z)^{-1}\|
 \le\frac1{|\operatorname{Im}z|-C_\omega}
 \quad\text{when }|\operatorname{Im}z|>C_\omega .
 \tag{20}
\] The resolvent identity at \(z=i(1+2C_\omega)\) gives an actual
Hilbert--Schmidt resolvent.

The general theorem in Mityagin--Siegl {[}MS, Theorem 2.2{]} says that a
closed densely defined operator with one \(\mathcal S_p\) resolvent has
complete generalized root spaces if its resolvent has polynomial bounds
on rays with successive gaps strictly less than \(\pi/p\). Take \(p=2\)
and the eight nonreal rays \(\pi/8+j\pi/4\), \(j=0,\ldots,7\). Equation
(20) gives the required bounds at large radius; the gap is
\(\pi/4<\pi/2\). Thus the hypotheses of the completeness theorem hold.
No Riesz-basis property is needed.

The same argument applies to
\(\mathcal A(\omega)^*=\mathcal A(\overline\omega)\). Therefore, if
every finite Riesz projection annihilates a vector, that vector is zero,
because it is orthogonal to the complete generalized root spaces of the
adjoint. This detection statement replaces any presumed convergent
expansion in ordinary eigenvectors.

\section{Full PDE reduction at defective angular
frequencies}\label{full-pde-reduction-at-defective-angular-frequencies}

Let \(P_\lambda^+\) be a Riesz projection in the \(i\) chiral sector. On
the two-sector direct sum use \[
 \mathcal P_\lambda=
 \operatorname{diag}(P_\lambda^+,\mathcal I P_\lambda^+\mathcal I).
 \tag{21}
\] It commutes with \(\mathcal A,\mathcal I,\mathcal J\) by their
resolvent identities. This refined pairing avoids double counting if
both \(\lambda,-\lambda\) occur in one sector and works at
\(\lambda=0\). Adjoint-root completeness detects any nonzero full PDE
solution at some interior radius with such a nonzero finite projection.

Choose Jordan chains \((\mathcal A-\lambda)u_j=u_{j-1}\), \(u_0=0\), in
the \(i\) sector; set \(v_j=\mathcal Iu_j\). Then \[
 \mathcal Av_j=-\lambda v_j-v_{j-1},\quad
 \mathcal Iu_j=v_j,\quad \mathcal Iv_j=u_j,\quad
 \mathcal Ju_j=iv_j,\quad \mathcal Jv_j=-iu_j .
 \tag{22}
\] The two chiral chains remain independent when \(\lambda=0\). If their
radial coefficients are \(p_j,q_j\), equation (2) reads \[
 \sqrt\Delta\,p_j'=\lambda p_j+p_{j+1}+(R_0-iR_1)q_j,\qquad
 \sqrt\Delta\,q_j'=-\lambda q_j-q_{j+1}+(R_0+iR_1)p_j .
 \tag{23}
\] Choose a chain and its highest index whose pair is not identically
zero. Higher coefficients vanish identically, leaving a homogeneous
ordinary two-component radial equation. The coefficient map treats
\(u_j,\mathcal Iu_j\) in the same way and intertwines
\(\mathcal I,\mathcal J\). The finite projections and coefficient maps
are bounded and independent of radius. They preserve the different
horizon weights in (7) and the infinity bounds.

Equivalently, use the finite invariant filtration
\(\ker(B-\lambda)^j\oplus\mathcal I\ker(B-\lambda)^j\). The radial
equation preserves it, and the first nonzero quotient has semisimple
angular action. This version does not apply an anticommuting angular
operator to the entire PDE.

In the radial basis \(Z=z_1(u+\mathcal Iu)-iz_2(u-\mathcal Iu)\), the
pair is \[
 (\partial_r\mp\Delta^{-1/2}R_0)z_{1/2}
 =\Delta^{-1/2}(R_1\mp i\lambda)z_{2/1}.
 \tag{24}
\] The two horizon powers are \((-1/4,-3/4)\) for positive spin,
\((3/4,5/4)\) for negative spin, multiplied by \(H\). It suffices to
exclude (24) for every genuine angular eigenpair. No uniform estimate in
\(m,\omega,\lambda\) has been assumed.

\section{A radial equation without scalar
division}\label{a-radial-equation-without-scalar-division}

Define \[
\begin{array}{c|cc}
 &w&w_\sharp\\ \hline
 \sigma=+1&
 \Delta^{3/4}H^{-1}e^{i\eta\omega r}z_2&
 \Delta^{1/4}H^{-1}e^{i\eta\omega r}z_1\\
 \sigma=-1&
 \Delta^{-3/4}H^{-1}e^{i\eta\omega r}z_1&
 \Delta^{-5/4}H^{-1}e^{i\eta\omega r}z_2 .
\end{array}
 \tag{25}
\] Both functions are smooth at the horizon, and \[
 w'=A w_\sharp,\qquad
 \Delta w_\sharp'+
 \left(\frac{1-3\sigma}{2}\Delta'+2i\widetilde K\right)w_\sharp=B w,
 \tag{26}
\] where \(A=R_1+i\sigma\lambda\), \(B=R_1-i\sigma\lambda\). Put \[
 p=\Delta^{-(1+3\sigma)/2}H^2e^{-2i\eta\omega r},\qquad
 v_r=pw,\qquad z_r=\Delta p w_\sharp .
 \tag{27}
\] The subscripts distinguish these radial functions from angular
expectations. Direct differentiation gives \[
 (\Delta\partial_r+L_r)v_r=A z_r,\qquad z_r'=Bv_r,\qquad
 L_r=\frac{1+3\sigma}{2}\Delta'-2i\widetilde K .
 \tag{28}
\] Since \(A'=B'=R_1'\) is constant, two derivatives of the first
equation, using the second, give \[
 \mathcal E v_r=0,\qquad
 \mathcal E=
 \partial_r^2(\Delta\partial_r+L_r)
 -(R_1^2+\lambda^2)\partial_r-(3R_1-i\sigma\lambda)R_1' .
 \tag{29}
\] All coefficients are polynomials of degree at most two in \(r\).
Neither \(A^{-1}\) nor \(B^{-1}\) is used. Apparent scalar poles
therefore cause no loss of solutions.

\section{Distributional horizon
extension}\label{distributional-horizon-extension}

Write \(t=r-r_+\), \(d_H=\Delta'(r_+)=2\delta\). Locally \[
 v_r=t^\alpha f(t),\quad z_r=t^{\alpha+1}g(t),\quad f,g\in C^\infty,
 \qquad
 \alpha=-\frac{1+3\sigma}{2}+\frac{2iK(r_+)}{d_H}.
 \tag{30}
\] Negative spin has \(\operatorname{Re}\alpha>1\), positive spin has
\(\operatorname{Re}\alpha>-2\), since
\(\operatorname{Re}(2iK(r_+)/d_H)=2(r_+^2+a^2)\gamma/d_H\).

\subsection{Generic frequencies}\label{generic-frequencies}

For negative spin ordinary zero extension solves (28) distributionally.
For positive spin with \(\alpha\ne-1\), use the canonical meromorphic
distributions \(t_+^\alpha\) and \(t_+^{\alpha+1}\), multiplied by
\(f,g\), and joined to the original solution away from the horizon.
Their only pole in the stated region is \(\alpha=-1\). The identities \[
 t\partial_t(t_+^\alpha f)
 =\alpha t_+^\alpha f+t_+^{\alpha+1}f',\qquad
 \partial_t(t_+^{\alpha+1}g)
 =(\alpha+1)t_+^\alpha g+t_+^{\alpha+1}g'
 \tag{31}
\] hold by meromorphic continuation from the integrable region,
including at nonnegative integers. With \(\Delta=t(d_H+t)\), the
residuals in (28) are the canonical extensions of identically vanishing
smooth coefficient residuals. Thus the extended pair is homogeneous.
This is an argument at fixed \(\omega\), requiring no holomorphic family
of physical outgoing solutions.

\subsection{The exceptional positive-spin
resonance}\label{the-exceptional-positive-spin-resonance}

The only exceptional negative integer is \(\alpha=-1\), at \[
 \omega=\omega_H=\frac{am+id_H/2}{B_+},\qquad
 B_+=r_+^2+a^2,\qquad \gamma=\frac{\delta}{B_+}.
 \tag{32}
\] At this point \(L_r(r_+)=d_H\). Define \(G_H=2Q(Mr_+-Q^2)/d\). If
\(B(r_+)=0\), then \(\lambda=\beta+iG_H\omega\), so
\(\lambda^2-\beta^2=2i\beta G_H\omega-G_H^2\omega^2\). Compare with
(18), divide by \(\omega\ne0\), and take imaginary parts: \[
 0=-\frac{\gamma T}{|\omega|^2}
       +\gamma(D+G_H^2)-2\beta G_H .
 \tag{33}
\] But \(T\ge0\), \(D\le a^2\kappa^2\), and \[
 G_H^2+a^2\kappa^2=a^2+\frac{4Q^2B_+}{d},
\] so \[
 2\beta G_H-\gamma(G_H^2+a^2\kappa^2)
 =6M+\delta\left(\frac{6M^2-4Q^2}{d}-\frac{a^2}{B_+}\right)>0.
 \tag{34}
\] The first quotient in parentheses is \(4+2M^2/d>6\), while
\(a^2/B_+<1\). Both signs of \(Q\) are included. This contradicts (33),
so \(B(r_+)\ne0\). The same squared argument would exclude \(A(r_+)=0\),
but that is unnecessary below.

Here \(z_r\) is smooth and \(v_r=t^{-1}f\). The equation \(z_r'=Bv_r\)
forces \(f(0)=0\), so \(v_r\) is also smooth. Its zero extension need
not be homogeneous: \(z_r'\) acquires a delta source. The required
corrected extension is \[
 z_e=1_{[r_+,\infty)}z_r,\qquad
 v_e=1_{[r_+,\infty)}v_r+
       \frac{z_r(r_+)}{B(r_+)}\delta_{r_+}.
 \tag{35}
\] Since \[
 (\Delta\partial_r+L_r)\delta_{r_+}
 =(-d_H+L_r(r_+))\delta_{r_+}=0,
\] the first equation of (28) is unchanged. The added term in \(Bv_e\)
precisely matches the delta coefficient in \(z_e'\). Thus (28), and
hence (29), hold homogeneously.

This correction is essential: vanishing of the Laurent coefficient alone
does not remove a finite horizon boundary source. The corrected
extension is of finite distribution order, supported in
\([r_+,\infty)\), and agrees with the physical function on the open
exterior.

\section{Complex Laplace transform and positive
energy}\label{complex-laplace-transform-and-positive-energy}

\subsection{Transform equation}\label{transform-equation}

At infinity the homogeneous extension \(v_e\) agrees with
\(e^{-i\eta\omega r}O_{\rm con}(r^C)\), by (25)--(27). Its
Fourier--Laplace transform \[
 F(\xi)=\langle v_e(r),e^{-i\xi r}\rangle
 \tag{36}
\] is holomorphic for \(\operatorname{Im}\xi<-\eta\gamma\). The pairing
is defined by exponential damping and cutoffs; compactly supported
finite-part or delta terms present no convergence problem. Restrict to
the interior ray \[
 \xi=\omega x,\qquad -\infty<x<-h,\qquad Y(x)=F(\omega x).
 \tag{37}
\] Here \(\gamma(x+\eta)<-2\gamma\).

Transform (29) by \(\partial_r\mapsto i\xi\), \(r\mapsto i\partial_\xi\)
in the correct operator order, then set \(\xi=\omega x\). The result is
\[
 a_xY''+(b_x+i\omega c_x)Y'
 +(d_0+\omega d_1+\omega^2d_2+i\omega e_x)Y=0,
 \tag{38}
\] where \[
\begin{aligned}
 a_x&=x(x+\eta)(x+h),&
 b_x&=x^2-\eta h-3\sigma x(x+h),\\
 c_x&=2Mx(x+\eta)^2,&
 e_x&=M(x+\eta)(x-\eta-3\sigma x),\\
 d_0&=(\beta^2-\lambda^2)x+\tfrac32\sigma h-c\lambda,&
 d_1&=2amx^2,\\
 d_2&=-[2a^2+(a^2+Q^2)(x+2\eta)]x^2-Q^2\eta^2x .
\end{aligned}
 \tag{39}
\] The coefficients in (39) follow by applying the distributional
Fourier substitution to the polynomial radial operator (29), preserving
the order of multiplication and differentiation.

Set \[
 \mu=x^{-2}(x+\eta)^{-3\sigma},\quad
 P=\mu a_x=\frac{x+h}{x(x+\eta)^{3\sigma-1}},\quad
 W_x=\frac{(x+\eta)^{-3\sigma}}x .
 \tag{40}
\] Both \(P,W_x\) are strictly positive on this interval. With
\(q_x=\mu c_x/2=M(x+\eta)^{2-3\sigma}/x\), \[
 (\mu a_x)'=\mu b_x,\qquad q_x'=\mu e_x,\qquad
 q_x/P=\frac{M(x+\eta)}{x+h}.
 \tag{41}
\] Define \[
 S(x)=M[x-2\log(-x-h)],\qquad Y=e^{-i\omega S}U .
 \tag{42}
\] Because \(S'=q_x/P\), the imaginary transport disappears: \[
 (PU')'+W_x\left[
 \beta^2-\lambda^2+\frac{3\sigma h/2-c\lambda}{x}
 +2am\omega x+\omega^2R(x)\right]U=0,
 \tag{43}
\] \[
 R(x)=(x+\eta)^2\left(\delta^2-\frac{2M^2}{x+h}\right)
      +a^2\eta^2-2a^2x .
 \tag{44}
\]

\subsection{The coercive sign}\label{the-coercive-sign}

Use the genuine angular eigenvector in (18). The bracket in (43) equals
\[
 -k_\sigma(x)+\omega l(x)+\omega^2\mathcal B(x),
 \tag{45}
\] where \[
 k_\sigma=T+\frac c x(L-\sigma\beta)
       =\langle K_\sigma(x)\rangle\ge0,\qquad
 l(x)=2amx-C-\frac c x v\in\mathbb R,\qquad
 \mathcal B=R-D.
\] For \(y=-x-\eta>2\), equation (19) gives the exact identity \[
 \mathcal B(x)=
 y^2\left(\delta^2+\frac{2M^2}{y-2}\right)
 +2a^2y-a^2\langle\sin^2\theta\rangle>0 .
 \tag{46}
\] These signs hold for every allowed background and actual complex
angular eigenpair. No complex separation constant is treated as an
arbitrary real parameter. The real, possibly indefinite coefficient
\(l\) will cancel in the imaginary-part identity.

\subsection{Endpoint justification}\label{endpoint-justification}

At \(x\uparrow-h\), the Laplace transform is analytic across the
endpoint, which is strictly inside its holomorphy half-plane. Writing
\(t_x=-x-h\), equation (42) gives \[
 |U|=O(t_x^{2M\gamma}),\qquad
 |U'|=O(t_x^{2M\gamma-1}).
 \tag{47}
\] Here \(P\asymp t_x\), \(W_x\asymp1\), \(\mathcal B\asymp t_x^{-1}\),
and \(k_\sigma,l\) are bounded. The worst energy integrands behave as
\(t_x^{4M\gamma-1}\), which is integrable for every \(\gamma>0\). The
flux \(PU'\overline U\) tends to zero as \(O(t_x^{4M\gamma})\).

At \(x\to-\infty\), an explicit finite-part representation gives a
polynomial bound times \(e^{\gamma r_+x}\) for \(F(\omega x)\) and each
fixed derivative. Cut off near the horizon and absorb the cutoff into
the smooth \(f\). For positive spin at a generic frequency, use \[
 t_+^\alpha f=
 \frac{\partial_t(t_+^{\alpha+1}f)}{\alpha+1}
 -\frac{t_+^{\alpha+1}f'}{\alpha+1}.
 \tag{48a}
\] Both undifferentiated distributions on the right are integrable since
\(\operatorname{Re}\alpha>-2\). Their absolute Laplace integrals are
bounded by \[
 e^{\gamma r_+x}
 \int_0^\varepsilon t^{\operatorname{Re}\alpha+1}
             e^{\gamma xt}\,dt
 =e^{\gamma r_+x}O(|x|^{-\operatorname{Re}\alpha-2}).
\] The distributional derivative contributes at most one factor of
\(|x|\). Negative spin is integrable directly; the exceptional extension
is an ordinary smooth zero extension plus a delta. Inserting \(r^j\)
handles each derivative in \(x\) by the same argument. The smooth
outgoing tail begins a positive distance from the horizon and has a
stronger exponential bound for sufficiently large negative \(x\). Thus
no general support-edge distribution bound or Watson asymptotic
coefficient is being presumed.

The gauge factor has modulus \(e^{-\gamma Mx}\) times a power of
\(-x-h\). Thus, after adjusting constants, \[
 |U|+|U'|\le C(1+|x|)^N e^{\gamma(r_+-M)x}
          =C(1+|x|)^N e^{\gamma\delta x}.
 \tag{48}
\] Because \(\delta>0\), this decays exponentially. All weights have at
most polynomial growth at this endpoint. The energy integrals converge
and the boundary flux is zero, including for the delta correction (35).

\subsection{Energy identity and
injectivity}\label{energy-identity-and-injectivity}

Multiply (43) by \(\overline U\), integrate, and use the established
endpoint bounds: \[
 -\int P|U'|^2-\int W_x k_\sigma|U|^2
 +\omega\int W_xl|U|^2
 +\omega^2\int W_x\mathcal B|U|^2=0 .
\] Multiplication by \(\overline\omega\) and taking imaginary parts
eliminates the real linear coefficient: \[
 0=\gamma\left[
  \int P|U'|^2+\int W_x k_\sigma|U|^2
  +|\omega|^2\int W_x\mathcal B|U|^2\right].
 \tag{49}
\] Since \(\gamma>0\), \(\omega\ne0\), \(P,W_x>0\), \(k_\sigma\ge0\),
and \(\mathcal B>0\), this forces \(U=0\).

The transform \(F\) vanishes on an interior complex ray, hence on its
entire half-plane by the identity theorem. On any sufficiently low
horizontal line it is the Fourier transform of an exponentially damped
tempered distribution. Fourier injectivity gives \(v_e=0\), in
particular \(v_r=0\) in the open exterior. Since \(p\ne0\), \(w=0\);
(26) then gives \(Aw_\sharp=0\). The affine polynomial \(A\) is not
identically zero because \(Q\omega M/d\ne0\). Its possible zero is
isolated, so continuity gives \(w_\sharp=0\) everywhere. The invertible
transformations give \(z_1=z_2=0\).

This contradicts the nonzero pair extracted from any nonzero full PDE
solution. Both systems (2) vanish, and their invertible source
transformations imply the radiative conclusion (8a)--(8b) for \(Q\ne0\).

\section{The zero-charge Kerr case}\label{the-zero-charge-kerr-case}

For \(Q=0\), the background Maxwell field vanishes. Linearization then
separates the gravitational and Maxwell radiative equations into the
standard Teukolsky equations of spins \(s\in\{-2,-1,1,2\}\). The Kerr
mode-stability input used below is Whiting's upper-half-plane theorem
{[}W{]} in the separated form recalled in {[}AMPW, Introduction and
Section II.A.3, Definition 3{]}, restricted to negative spins
\(s=-1,-2\): a regular-axis separated solution which is the
\emph{non-special ingoing Frobenius solution} at the horizon and
outgoing at infinity must vanish. Positive spins will be reduced to
negative spins by the spin-reversal lemma below. A general
\(e^{-i\omega t+im\phi}F(r,\theta)\) need not initially be a product, so
the full-field reduction is also proved here.

On the fixed-\(m\) subspace of \(L^2((0,\pi),\sin\theta\,d\theta)\),
define the positive angular convention \[
 T_s(\omega)=
 -\frac1{\sin\theta}\partial_\theta(\sin\theta\,\partial_\theta)
 +\frac{(m+s\cos\theta)^2}{\sin^2\theta}
 -a^2\omega^2\cos^2\theta+2as\omega\cos\theta .
 \tag{50}
\] The zero-frequency realization \(T_s(0)\) is the canonical
self-adjoint spin-bundle operator restricted to this azimuthal subspace.
Its orthonormal spin-weighted spherical harmonics have eigenvalues \[
 \Lambda_\ell=\ell(\ell+1)-s^2,\qquad
 \ell\ge\max(|m|,|s|).
 \tag{51}
\] The domain imposes the regular spin-weighted branch. This distinction
is essential when \(m\pm s=0\): a logarithmic scalar solution is also in
\(L^2(\sin\theta\,d\theta)\), but it is excluded by the
canonical/Friedrichs domain. We do not use the maximal scalar interval
domain. Smooth eigenfunctions and generalized eigenfunctions follow from
elliptic regularity on the spin bundle after addition of the smooth
potential.

For fixed \(\omega\), the last two terms of (50) are a bounded complex
multiplication operator. Since (51) grows quadratically in \(\ell\), a
fixed-\(m\) resolvent of \(T_s(0)\) is trace class. The same
bounded-perturbation and nonreal-ray argument used in (20), now with
\(p=1\), proves generalized-root completeness for \(T_s(\omega)\) and
its adjoint. For example the same eight rays have gaps \(\pi/4<\pi\).

For clarity, the canonical eigenvalue \(\Lambda\) in (50) is not the
standard radial separation constant. With
\(K_{\rm std}=(r^2+a^2)\omega-am\), direct separation of the physical
Teukolsky equation gives \[
\begin{split}
0={}&\Delta^{-s}\partial_r(\Delta^{s+1}\partial_r R)
 +\frac{K_{\rm std}^2-2is(r-M)K_{\rm std}}{\Delta}R\\
 &+(4is\omega r+s-\Lambda-a^2\omega^2+2am\omega)R .
\end{split}
 \tag{52}
\] Thus the standard radial constant is
\(\Lambda-s+a^2\omega^2-2am\omega\). The fixed-frequency PDE is a sum of
a radial operator and \(-T_s(\omega)\), with precisely this scalar
shift.

\subsection{Horizon branches and the spin-reversal
identity}\label{horizon-branches-and-the-spin-reversal-identity}

For the Kerr background set \[
 \Omega_H=\frac{a}{r_+^2+a^2},\qquad
 \kappa_H=\frac{r_+-r_-}{2(r_+^2+a^2)},\qquad
 k=\omega-m\Omega_H.
\] In this section \(H'/H=-iK_{\rm std}/\Delta\). Writing
\(R=H\Delta^{-s}f\), direct conjugation of (52) gives \[
 E_sf=0,\qquad E_s=\Delta\partial_r^2+P_s\partial_r+C_s,
 \tag{53}
\] with \[
 \begin{aligned}
 P_s&=(1-s)\Delta'-2iK_{\rm std},\\
 C_s&=2i(2s-1)\omega r-s-\Lambda-a^2\omega^2+2am\omega .
 \end{aligned}
\] Indeed, for \(\psi=\partial_r\log(H\Delta^{-s})
=-(iK_{\rm std}+s\Delta')/\Delta\), the first-derivative coefficient
after conjugation is \(2\Delta\psi+(s+1)\Delta'=P_s\). The remaining
\(1/\Delta\) terms cancel against the potential in (52), using
\(\Delta'=2(r-M)\) and \(K_{\rm std}'=2\omega r\), leaving the displayed
\(C_s\).

\textbf{Lemma 8.1 (negative-spin horizon branch).} Let \(s\in\{-1,-2\}\)
and \(\Lambda\in\mathbb C\). If a nonzero solution of (52) satisfies
\(R\in H\Delta^{-s}C^\infty([r_+,r_++\varepsilon))\), then \[
 R=c\,t^{-s-ik/(2\kappa_H)}(1+O(t)),\qquad
 t=r-r_+,\quad c\ne0.
 \tag{54}
\] In particular, it is the non-special ingoing Frobenius solution
required by {[}AMPW, Definition 3{]}.

\emph{Proof.} Write \(d_H=r_+-r_->0\), so \(\Delta=t(d_H+t)\). The
indicial roots of (53) for \(f\) are \[
 0,\qquad k_s=s+\frac{ik}{\kappa_H},\qquad
 \operatorname{Re}k_s=s-\frac{\gamma}{\kappa_H}<0.
\] The root \(0\) admits a unique analytic solution with a freely chosen
nonzero constant coefficient. The other independent solution has a
nonzero \(t^{k_s}\) contribution, possibly accompanied by logarithms at
integral root differences, and cannot be smooth at \(t=0\). A nonzero
smooth \(f\) must therefore be a nonzero multiple of the root-\(0\)
solution. Since \(H=t^{-ik/(2\kappa_H)}\) times a smooth nonvanishing
factor and \(\Delta^{-s}\sim d_H^{-s}t^{-s}\), (54) follows. \(\square\)

\textbf{Lemma 8.2 (spin reversal, including exceptional frequencies).}
Fix \(s\in\{1,2\}\), \(m\in\mathbb Z\), and
\(\operatorname{Im}\omega>0\). Let \(T_s(\omega)u=\Lambda u\) with
\(u\ne0\) in the canonical angular domain, and let \(R\ne0\) solve (52).
Suppose \(R\in H\Delta^{-s}C^\infty([r_+,\infty))\) and \(HR\) has
polynomial conormal bounds at infinity. Set \(f=H^{-1}\Delta^sR\), and
define \[
 \widetilde u(\theta)=u(\pi-\theta),\qquad
 \widetilde R=H\Delta^s\partial_r^{2s}f.
 \tag{55}
\] Then \(\widetilde u\) is a canonical spin-\((-s)\) angular
eigenfunction with eigenvalue \(\Lambda\), and \(\widetilde R\) is a
nonzero spin-\((-s)\) radial solution satisfying the same future-horizon
regularity and outgoing infinity conditions. This holds at every
frequency in the open upper half-plane, without a nonvanishing horizon
leading-coefficient assumption.

\emph{Proof.} Reflecting \(\theta\mapsto\pi-\theta\) in (50)
interchanges \(T_s(\omega)\) and \(T_{-s}(\omega)\) at the same
\((\omega,m,\Lambda)\). It maps the regular spin-bundle domains into one
another: the pole exponents \(|m+s|\) and \(|m-s|\) are exchanged. This
proves the angular claim, including its domain statement.

For the radial equation put \(n=2s\). Because \(\Delta\) and \(P_s\) are
quadratic and \(C_s\) is affine in \(r\), Leibniz's rule applied to (53)
yields \[
 \boxed{\quad\partial_r^{2s}E_s
 = E_{-s}\partial_r^{2s}.\quad}
 \tag{56}
\] To verify all coefficients, the coefficient of \(\partial_r^{n+1}\)
is \(n\Delta'+P_s=P_{-s}\); that of \(\partial_r^n\) is \[
 \frac{n(n-1)}2\Delta''+nP_s'+C_s=C_{-s};
\] and the coefficient of \(\partial_r^{n-1}\) vanishes since \[
 \frac{n(n-1)}2P_s''+nC_s'
 =2i\omega n(2s-n)=0.
\] There are no further lower-order terms. Thus \(g=\partial_r^{2s}f\)
satisfies \(E_{-s}g=0\), and the conjugation (53) for \(-s\) proves that
\(\widetilde R=H\Delta^sg\) solves the spin-\((-s)\) version of (52).
Since \(f\) is smooth at the horizon, so is \(g\). This establishes
\(\widetilde R\in H\Delta^sC^\infty\), including the resonant Frobenius
cases.

It remains to exclude a zero image and check infinity. For \(\gamma>0\),
the two radial exponential types are distinct, and polynomial conormal
control of \(HR\) selects the outgoing solution. The outgoing asymptotic
expansion of (52), including differentiated expansions, gives \[
 f=e^{2i\omega r}r^{4iM\omega-1}
       \bigl(c_0+O(r^{-1})\bigr),\qquad c_0\ne0.
 \tag{57}
\] Here \(H=e^{-i\omega r}r^{-2iM\omega}
(\text{constant}+O(r^{-1}))\). Consequently \[
 \partial_r^{2s}f=(2i\omega)^{2s}
 e^{2i\omega r}r^{4iM\omega-1}
       \bigl(c_0+O(r^{-1})\bigr).
 \tag{58}
\] Since \(\omega\ne0\), this is not identically zero. Furthermore
\(H\widetilde R=H^2\Delta^s\partial_r^{2s}f\) has polynomial conormal
bounds and belongs to the spin-\((-s)\) outgoing branch. Thus
\(\widetilde R\) is nonzero, with the stated boundary conditions.
\(\square\)

For an angular eigenfunction the usual Kerr admissibility inequality is
also automatic in our normalization: \[
 \operatorname{Im}(\overline\omega\Lambda)\|u\|^2
 =-\gamma\left(\langle u,T_s(0)u\rangle
             +a^2|\omega|^2\|\cos\theta\,u\|^2\right)<0,
 \tag{59}
\] since the canonical \(T_s(0)\) has eigenvalues
\(\ell(\ell+1)-s^2\ge |s|>0\) for \(s=\pm1,\pm2\).

\subsection{Completion of the full Kerr field
argument}\label{completion-of-the-full-kerr-field-argument}

Any nonzero full PDE field has a nonzero finite Riesz projection by
adjoint-root completeness. On a Jordan chain
\((T_s-\Lambda)u_j=u_{j-1}\), its highest nonzero radial coefficient
solves (52) homogeneously. The associated bottom angular vector is a
genuine regular eigenfunction of (50). Projections and coefficient maps
are bounded and independent of \(r\), so they preserve the scalar
spin-dependent horizon factor \(H\Delta^{-s}\), the smooth horizon
coefficient and the outgoing infinity exponential with polynomial
bounds. At infinity, the two radial asymptotic exponentials are distinct
for \(\gamma>0\); the polynomial bound after multiplication by \(H\)
selects the outgoing one. The resulting nonzero separated mode has a
regular angular eigenfunction, weighted-smooth horizon data and outgoing
infinity behavior. If \(s=-1,-2\), Lemma 8.1 identifies its horizon
behavior with the non-special ingoing Frobenius branch, so Whiting's
negative-spin theorem excludes it. If \(s=1,2\), Lemma 8.2 produces a
nonzero separated outgoing mode of spin \(-s\) satisfying the same
negative-spin theorem. Either possibility is a contradiction. Descending
the finite Jordan chain and invoking adjoint-root completeness, the
entire fixed-frequency Kerr radiative field vanishes, including for
\(a=0\).

\textbf{Separation conventions.} The radial equation is the physical
Teukolsky equation {[}T{]}, and its symmetrized operator formulation is
also discussed in {[}PV{]}. The canonical angular eigenvalue \(\Lambda\)
in (50) differs by a scalar shift from the radial separation constants
common in the literature. The conversion between (50)--(52) and the AMPW
normalization \(\widetilde R=\Delta^{s/2}R\) is specified after {[}AMPW,
equation (7){]}:
\(\Lambda_{\rm AMPW}+2am\omega-a^2\omega^2=A_{\rm Whiting}+s\), where
\(A_{\rm Whiting}=\Lambda-s\) in our convention. Only the negative-spin,
non-special form of the Kerr theorem is invoked. In particular the
argument does not divide by a positive-spin horizon
Teukolsky--Starobinsky coefficient or require a positive-spin horizon
leading coefficient to be nonzero.

This completes the zero-charge sector, and hence the theorem.
\(\square\)

\begin{thebibliography}{AMPW}
\bibitem[H]{H} P. Hintz, ``Perturbations of Kerr--Newman black holes:
separation and mode stability,'' arXiv:2609.33661v1 (2026).
\url{https://arxiv.org/abs/2609.33661}.

\bibitem[MS]{MS} B. Mityagin and P. Siegl, ``Local form
subordination without a power decay and a criterion of Riesz
basesness,'' \emph{Bol. Soc. Mat. Mex.} \textbf{31}, 111 (2025),
Theorem 2.2. \url{https://doi.org/10.1007/s40590-025-00751-7}.

\bibitem[W]{W} B. F. Whiting, ``Mode stability of the Kerr black hole,''
\emph{J. Math. Phys.} \textbf{30}, 1301--1305 (1989).
\url{https://doi.org/10.1063/1.528308}.

\bibitem[AMPW]{AMPW} L. Andersson, S. Ma, C. Paganini, and B. F. Whiting,
``Mode stability on the real axis,'' \emph{J. Math. Phys.} \textbf{58},
072501 (2017). \url{https://arxiv.org/abs/1607.02759}.

\bibitem[TdC]{TdC} R. Teixeira da Costa, ``Mode stability for the
Teukolsky equation on extremal and subextremal Kerr spacetimes,''
\emph{Commun. Math. Phys.} \textbf{378}, 705--781 (2020).
\url{https://doi.org/10.1007/s00220-020-03796-z}.

\bibitem[PV]{PV} O. Petersen and A. Vasy, ``Dual modes in Kerr
spacetimes and the Whiting transform: Mode stability revisited,''
arXiv:2608.26034 (2026). \url{https://arxiv.org/abs/2608.26034}.

\bibitem[T]{T} S. A. Teukolsky, ``Perturbations of a rotating black
hole. I. Fundamental equations for gravitational, electromagnetic,
and neutrino-field perturbations,'' \emph{Astrophys. J.}
\textbf{185}, 635--647 (1973). \url{https://doi.org/10.1086/152444}.

\bibitem[He]{He} L. He, ``The linear stability of weakly charged and
slowly rotating Kerr--Newman family of charged black holes,''
\emph{Annals of PDE} \textbf{12}, 3 (2026).
\url{https://doi.org/10.1007/s40818-025-00219-x}.

\bibitem[DGS]{DGS} O. J. C. Dias, M. Godazgar, and J. E. Santos,
``Linear mode stability of the Kerr--Newman black hole and its
quasinormal modes,'' \emph{Phys. Rev. Lett.} \textbf{114}, 151101
(2015). \url{https://doi.org/10.1103/PhysRevLett.114.151101}.

\bibitem[DCH]{DCH} K. Destounis, V. Cardoso, and P. Hintz,
``The quasinormal modes of Kerr--Newman black holes: a separated approach,''
arXiv:2610.07142 (2026). \url{https://arxiv.org/abs/2610.07142}.

\bibitem[M]{M} H. Motohashi, ``Charge-regular formulation and radiative
amplitudes of Kerr--Newman perturbations,'' arXiv:2610.10529 (2026).
\url{https://arxiv.org/abs/2610.10529}.

\bibitem[Z]{Z} J.-K. Zhao, ``Exact algebraic separation of the
gravitoelectromagnetic perturbations of Kerr--Newman--(anti-)de Sitter
black holes,'' arXiv:2610.22274 (2026).
\url{https://arxiv.org/abs/2610.22274}.
\end{thebibliography}

\end{document}
